\documentclass{article}
\usepackage{array}

\title{Compos3D Profiling Report}
\author{}
\date{}

\begin{document}
\maketitle

\section{Performance Analysis}
\subsection{Profiling Execution time}
We profiled five important functions in the Compos3D codebase using Python's standard library \texttt{profile} module. The measurements summarized in this section are taken directly from \texttt{profiling/important\_functions/summary.json}. They were collected with the local mock LLM path and the heuristic critic so that the profiling run exercised the real orchestration logic without depending on AWS or Blender. Because \texttt{profile} is a pure-Python profiler, the recorded values include profiler overhead; the results are therefore most useful for comparing relative costs and identifying where optimization work is likely to matter most.

\begin{table}[h]
\centering
\small
\renewcommand{\arraystretch}{1.15}
\begin{tabular}{|>{\raggedright\arraybackslash}p{0.34\linewidth}|>{\raggedright\arraybackslash}p{0.20\linewidth}|r|r|r|}
\hline
Function & File & Calls & Total (s) & Cumulative (s) \\
\hline
\texttt{run\_vertical\_inference} & \texttt{engine.py} & 150 & 0.009680 & 0.970212 \\
\hline
\texttt{train\_vertical\_slice} & \texttt{engine.py} & 1 & 0.000049 & 0.258973 \\
\hline
\texttt{SceneHypothesisLoop.train} & \texttt{loop.py} & 1 & 0.004125 & 0.184729 \\
\hline
\rule{0pt}{3.1ex}\shortstack[l]{\texttt{MockSceneLLM.}\\\texttt{generate\_scene\_program}} & \texttt{scene\_llm.py} & 2500 & 0.031871 & 0.130528 \\
\hline
\texttt{\_build\_heuristic\_score} & \texttt{critic.py} & 5000 & 0.024567 & 0.059346 \\
\hline
\end{tabular}
\caption{Profiled functions and recorded execution times.}
\end{table}

\paragraph{\texttt{run\_vertical\_inference}}
This function was profiled because it is the main inference entry point and covers bank loading, hypothesis selection, scene generation, scoring, and artifact writing. It had the highest cumulative time at 0.970212 seconds, while its own direct CPU time remained small. Most of the cost comes from repeated JSON serialization and filesystem writes. The clearest improvement is to reduce small file writes by batching outputs, reusing directories, and skipping non-essential artifacts in fast inference mode.

\paragraph{\texttt{train\_vertical\_slice}}
This function was chosen because it is the top-level training entry point exposed through the CLI. Its total self time is negligible at 0.000049 seconds, while its cumulative time is 0.258973 seconds, showing that nearly all work happens in lower layers. This confirms that the wrapper is not the bottleneck. The best improvement here is simply to keep the orchestration path lean and avoid unnecessary config copying or serialization.

\paragraph{\texttt{SceneHypothesisLoop.train}}
This is the core hypothesis training loop, so it was profiled as the main learning-time hot path. Its cumulative time was 0.184729 seconds for the local mock run. The detailed profile shows that much of that time is spent saving programs, snapshots, traces, and manifests rather than in the loop logic itself. A useful optimization would be to buffer outputs in memory and flush them less often, especially snapshot and trace files.

\paragraph{\texttt{MockSceneLLM.generate\_scene\_program}}
This function was included because scene generation sits directly on the hot path for both training and inference. Across 2500 calls it accumulated 0.130528 seconds. The main costs are prompt asset extraction, normalization, and Pydantic model construction. The strongest improvements would be caching repeated prompt parsing, avoiding repeated supported-asset normalization, and deferring heavy validation until the final output boundary.

\paragraph{\texttt{\_build\_heuristic\_score}}
This function was selected because every generated scene is scored with it, so it directly affects evaluation throughput and training feedback. Across 5000 calls it accumulated 0.059346 seconds. Its work is relatively small per call, but frequent enough to matter. The best improvements would be to postpone rounding until serialization, reuse derived asset sets where possible, and use a lighter internal result structure before building the final Pydantic model.

Overall, the profile suggests that the largest current bottleneck in the mock local workflow is artifact I/O rather than pure computation. The most impactful near-term optimization would therefore be reducing JSON and filesystem overhead in the training and inference pipelines. After that, the next most valuable target is the repeated parsing and model construction work inside scene generation and heuristic scoring.

\end{document}
